\subsection{A one-point Duhamel criterion and quantitative obstruction localization}
\label{sec:onepointduhamel}
The bounded-penalty condition in Corollary~\ref{cor:boundedpenaltyRH}
was imposed at every real Fourier argument.  This spatial
quantifier is unnecessary.  Positivity of the heat kernel
lets a single observation point detect any negative
Laguerre value anywhere in the backward heat slab.
The reduction is rigorous, but it is an equivalence
criterion and not an estimate of the actual theta
penalty.

Throughout this subsection put $T=1/2$,
$U_n(s,x)=L_{n,\,1/2-s}(x)$ and
$k_n=(n+1)(2n+1)$.  Let $P_h$ denote the
one-dimensional Gaussian heat semigroup of
generator $\tfrac12\partial_x^2$ and put
\begin{align}
E_{n,C}(s,x)
 &= [k_nU_{n+1}(s,x)-CU_n(s,x)]_+,\nonumber\\
\mathcal E_{n,C}(T,x)
 &=\int_0^T(P_{T-s}E_{n,C}(s,\cdot))(x)\,ds,\nonumber\\
\mathcal N_n(T,x)
 &=\int_0^T(P_{T-s}[-U_n(s,\cdot)]_+)(x)\,ds.
\label{eq:onepointvariables}
\end{align}
The theta transform and each of its fixed-order
derivatives are bounded uniformly on the closed
heat-time slab.  In particular
\[
M_n:=\sup_{0\le s\le T,\ x\in\mathbb R}
|k_nU_{n+1}(s,x)|<\infty,
\quad
A_n:=\sup_{0\le s\le T,\ x\in\mathbb R}
|U_n(s,x)|<\infty.
\]

\begin{theorem}[The central-point criterion for RH]
\label{thm:onepointcriterion}
For any one fixed $x_0\in\mathbb R$, the
following conditions are equivalent:
\begin{enumerate}
\item The Riemann hypothesis.
\item For every integer $n\ge1$,
$\sup_{C\ge1}\mathcal E_{n,C}(T,x_0)<\infty$.
\item For every integer $n\ge1$,
$\mathcal E_{n,C}(T,x_0)=o(C)$
as $C\to\infty$.
\item For every integer $n\ge1$,
$\mathcal N_n(T,x_0)=0$.
\end{enumerate}
In particular $x_0=0$ suffices:
there is no need to assume a Duhamel bound at
every Fourier height.
\end{theorem}
\begin{proof}
The positive-part Lipschitz inequality,
already established in
Theorem~\ref{thm:quantpenalty}, gives
uniformly in $x$
\begin{equation}
C\mathcal N_n(T,x)-TM_n
\le\mathcal E_{n,C}(T,x)
\le C\mathcal N_n(T,x)+TM_n.
\label{eq:onepointsandwich}
\end{equation}
Thus (2), (3), and (4) are equivalent
for each $n$ at the \emph{single} chosen
point $x_0$.

The Gaussian kernel
$p_h(y)=(2\pi h)^{-1/2}\exp(-y^2/(2h))$
is strictly positive for every $h>0$
and real $y$.  If
$U_n(s_*,y_*)<0$ at a point
$0<s_*<T$, continuity supplies
a nonempty open space--time rectangle on
which $[-U_n]_+$ is strictly positive.
Its convolution at $x_0$, integrated over
this rectangle, is strictly positive.
Therefore $\mathcal N_n(T,x_0)=0$
implies $U_n(s,y)\ge0$ for all
$0<s<T$ and $y\in\mathbb R$.
Continuity extends the inequality to
$s=T$, giving
$L_{n,0}(y)\ge0$ at every order and
Fourier argument.  By the classical
generalized Laguerre--P\'olya
characterization for the real entire
Riemann $\Xi$ transform, this gives RH.

Conversely, RH and the
de Bruijn--Newman real-zero theorem
show that $H_t$ has only real zeros
throughout $0\le t\le1/2$.
Its generalized Laguerre inequalities
give $U_n\ge0$ on the slab, and hence
$\mathcal N_n(T,x_0)=0$.  This
proves all equivalences.
\end{proof}

There is also a reduction of
\emph{all orders} to one
nonnegative scalar penalty. Define
\[
b_n=1+T(A_n+M_n)>0,\qquad
S(C)=\sum_{n=1}^{\infty}
2^{-n}\frac{\mathcal E_{n,C}(T,0)}{b_n}.
\]
For each $C\ge1$, the sum is finite,
because $0\le\mathcal E_{n,C}(T,0)
\le T(M_n+CA_n)$ and hence
$S(C)\le 1+C$.

\begin{corollary}[A single scalar RH criterion]
\label{cor:scalarpenalty}
With the preceding normalizations,
\[
\boxed{\quad
\mathrm{RH}
\quad\Longleftrightarrow\quad
\sup_{C\ge1} S(C)<\infty .
\quad}
\]
\end{corollary}
\begin{proof}
If RH holds, $U_n\ge0$ and
$\mathcal E_{n,C}\le TM_n$,
so $S(C)\le\sum_n2^{-n}=1$
for every $C\ge1$.
If RH is false, the generalized
Laguerre characterization supplies an
order $n$ and a real $X$ for which
$U_n(T,X)<0$.  Theorem~\ref{thm:onepointcriterion}
and \eqref{eq:onepointsandwich}
give $\mathcal E_{n,C}(T,0)\to\infty$
linearly in $C$.  This one
positive summand forces
$S(C)\to\infty$.
\end{proof}

The last reformulation is compact,
but should not be mistaken for a
new unconditional inequality:
both criteria still encode the
all-order real-zero problem.

\begin{theorem}[Quantitative detection of a negative
Fourier-height witness]\label{thm:negativewitness}
Fix $n\ge1$. Assume that for some
$X\in\mathbb R$ and $\delta>0$,
\[
U_n(T,X)\le-\delta.
\]
Let finite constants $A,B\ge0$
satisfy
\[
|\partial_sU_n(s,x)|\le A,\qquad
|\partial_xU_n(s,x)|\le B
\]
for $0\le s\le T$ and $x\in\mathbb R$.
Set
\[
\tau=\min\left(\frac T4,\,
\frac{\delta}{8(1+A)}\right),\quad
r=\min\left(1,\frac{\delta}{4(1+B)}\right),
\quad Y=|X|+r.
\]
Then $\tau,r>0$ and
\begin{equation}
\boxed{
\mathcal N_n(T,0)
\ge
\frac{\delta r\sqrt{\tau}}{2\sqrt\pi}
\exp\!\left(-\frac{Y^2}{2\tau}\right)>0.
}
\label{eq:negativewitnessmass}
\end{equation}
Consequently every $C>0$ obeys
\begin{equation}
\mathcal E_{n,C}(T,0)\ge
\frac{C\delta r\sqrt{\tau}}{2\sqrt\pi}
e^{-Y^2/(2\tau)}-TM_n.
\label{eq:negativewitnesspenalty}
\end{equation}
\end{theorem}
\begin{proof}
On the rectangle
\[
T-2\tau\le s\le T-\tau,\qquad
|y-X|\le r,
\]
the derivative bounds and the
definition of $\tau,r$ imply
\[
U_n(s,y)\le
-\delta+2\tau A+rB\le-\delta/2.
\]
There $h=T-s\in[\tau,2\tau]$ and
$|y|\le Y$, so
\[
p_h(y)\ge
\frac1{\sqrt{4\pi\tau}}
\exp\!\left(-\frac{Y^2}{2\tau}\right).
\]
Integrate the lower bound
$[-U_n(s,y)]_+\ge\delta/2$
over the rectangle of
time-length $\tau$ and
spatial-length $2r$ to obtain
\eqref{eq:negativewitnessmass}.
The penalty estimate follows from
\eqref{eq:onepointsandwich}.
\end{proof}

\paragraph{What the theta correlations do and do not provide.}
The constants $A,B,M_n$ are finite
unconditionally.  They can be bounded
by explicit absolute moments of the
positive theta kernel, because every
$x$-derivative of $H_t$ is an integral
against a power of $u$ and because
$\partial_sU_n$ is given by the
Laguerre heat hierarchy.
This makes \eqref{eq:negativewitnesspenalty}
an effective \emph{conditional} inequality.

However, an off-real zero at very large
Fourier height can generate a signal at
$x=0$ suppressed by
$e^{-Y^2/(2\tau)}$.  An absolute
numerical upper bound on the central
penalty, even an extremely small one,
does not establish its boundedness
as $C\to\infty$.  To prove RH by this
route one must derive from the
complete square-lattice correlation
identities that
\[
\sup_{C\ge1}\mathcal E_{n,C}(T,0)<\infty
\qquad\text{for every }n\ge1,
\]
\emph{without} assuming $U_n\ge0$.
We have not proved this inequality.
The one-point reduction,
scalar packaging and witness lower bound
are rigorous structural results,
not a full proof of the Riemann hypothesis.
